%
% Halaman Abstrak
%
% @author  Andreas Febrian
% @version 2.2.0
% @edit by Ichlasul Affan
%

\chapter*{Abstrak}
\singlespacing

\noindent \begin{tabular}{l l p{10cm}}
	\ifx\blank\npmDua
		Nama&: & \penulisSatu \\
		Program Studi&: & \programSatu \\
	\else
		Nama Penulis 1 / Program Studi&: & \penulisSatu~/ \programSatu\\
		Nama Penulis 2 / Program Studi&: & \penulisDua~/ \programDua\\
	\fi
	\ifx\blank\npmTiga\else
		Nama Penulis 3 / Program Studi&: & \penulisTiga~/ \programTiga\\
	\fi
	Judul&: & \judul \\
	Pembimbing&: & \pembimbingSatu \\
	\ifx\blank\pembimbingDua
    \else
        \ &\ & \pembimbingDua \\
    \fi
    \ifx\blank\pembimbingTiga
    \else
    	\ &\ & \pembimbingTiga \\
    \fi
\end{tabular} \\

\vspace*{0.5cm}

\noindent Algoritma-algoritma kriptografi pasca-kuantum menjadi semakin relevan ketika sistem keamanan digital yang umum digunakan saat ini ternyata rawan diserang oleh komputer kuantum. Dari sekian algoritma-algoritma kriptografi pasca-kuantum yang tersedia, ada algoritma enkapsulasi kunci bernama ML-KEM dan algoritma tanda tangan digital bernama ML-DSA yang dirilis oleh NIST (National Institute of Standards and Technology) pada tahun 2024. Memahami kedua algoritma tersebut bersama dengan dua algoritma reduksi kisi LLL dan BKZ yang digunakan untuk mengevaluasi tingkat keamanan ML-KEM dan ML-DSA sangat penting baik untuk studi teoritis dan penerapan praktis dari kriptografi pasca-kuantum. Walaupun keempat algoritma tersebut relevan, hanya sedikit sarana pembelajaran yang dapat menjelaskan cara kerja keempat algoritma tersebut dengan penjelasan yang mudah dipahami. Skripsi ini membahas pengembangan Latviz, sebuah alat visualisasi digital yang didesain untuk membantu para pembelajar memahami proses keseluruhan dari algoritma-algoritma berbasis kisi seperti ML-KEM, ML-DSA, LLL, dan BKZ. Latviz menyediakan penjabaran visual bertahap yang interaktif beserta penjelasan tambahan untuk fungsi-fungsi dan variabel-variabel dari keempat algoritma tersebut. Alat ini dievaluasi kinerjanya melalui sebuah survei terhadap 22 mahasiswa dengan tingkat pengetahuan kriptografi yang beragam, berfokus pada seberapa efektif Latviz dalam meningkatkan pemahaman mereka terhadap ML-KEM, ML-DSA, LLL, dan BKZ. Hasil dari survei tersebut menunjukkan bahwa para pengguna Latviz umumnya menilai Latviz sebagai sarana pembelajaran yang efektif dalam meningkatkan pemahaman mereka terhadap keempat algoritma tersebut, dengan masing-masing sebesar 77,3\%, 81,8\%, 72,7\%, dan 77,3\% responden menganggap Latviz sebagai alat pembelajaran yang efektif untuk ML-KEM, ML-DSA, LLL, dan BKZ. Secara keseluruhan, hasil dari survei tersebut menunjukkan bahwa Latviz dapat menjadi sarana pembelajaran yang efektif untuk membantu para pembelajar memahami algoritma ML-KEM, ML-DSA, LLL, dan BKZ. \\

\vspace*{0.2cm}

\noindent Kata kunci: \\ Kriptografi, Kriptografi Pasca-Kuantum, Visualisasi Algoritma, Visualisasi, ML-KEM, ML-DSA, LLL, BKZ, Kriptografi berbasis kisi \\

\setstretch{1.4}
\newpage
